%! Representing a Quantitative Variable with Graphs — Unit 1, Lesson 1.3 — AP Practice (Answer Key)
% EFFL, AP Statistics variant: AP-format practice, assigned but NOT scored — the cover's score
% column reads NA for this component. The graded practice is the `homework` component that
% follows it at the back of the packet.
% Free-response (a) spirals to Lesson 1.1 (classify the variable) and (d) to Lesson 1.2 (a bar
% graph is for categories) — the two lessons this one is built on top of.
% GENERATED FROM homework/main.tex. Each \writelines{n} in the blank is answered by an
% \ansline{} terse enough not to wrap past n lines — that keeps the page counts equal.
% NO teachernote here — scoring guidance lives in the lesson plan.
\documentclass[10pt]{article}
\usepackage{apstats-article}
\usepackage{apstats-key}   % pulls in -boxes; do NOT also load -boxes

\begin{document}

\pageheader{Unit 1, Lesson 1.3}{AP Practice --- Answer Key}
% NAMESTRIP: no \namedateperiod here — mirrors the blank.

\vspace{0.15in}

% Authored BYTE-IDENTICALLY in ap_practice/ and ap_practice_key/ — this box is what keeps the
% blank and the key the same number of pages.
\boxguard[8]
\begin{remindbox}
\small
This page is \textbf{not required and not scored} --- your graded homework is the last section
of this packet. These are AP-format reps: work them for the practice, and bring what you get
stuck on to office hours or study hall.
\end{remindbox}

\vspace{0.12in}

\boxguard[24]
\begin{scenariobox}[The Rainfall Record]{navy}
\small
A weather station recorded the total rainfall, in millimetres, on each of 20 rainy days. The
stem is the tens digit and each leaf is a ones digit, so the row \textbf{2} stands for
20, 22, 25, and 29.

\vspace{8pt}
\begin{center}
\renewcommand{\arraystretch}{1.25}
\begin{tabular}{@{} r | l @{}}
\textbf{0} & 2 \quad 4 \quad 5 \quad 8 \quad 9 \\
\textbf{1} & 1 \quad 3 \quad 3 \quad 6 \quad 7 \quad 8 \\
\textbf{2} & 0 \quad 2 \quad 5 \quad 9 \\
\textbf{3} & 1 \quad 4 \quad 8 \\
\textbf{4} & 2 \quad 7 \\
\end{tabular}
\end{center}
\end{scenariobox}

\vspace{0.14in}

\boxguard[16]
\begin{headlinebox}{goldbg}
\small
\textbf{Section I --- Multiple Choice.} Circle the single best answer. Every answer choice is
about the context, not just the arithmetic --- read them all before choosing.
\end{headlinebox}

\vspace{0.10in}

\begin{enumerate}[leftmargin=1.9em, itemsep=0.13in]

  \item On how many of these days did more than 30 mm of rain fall?
        \begin{enumerate}[label=\textbf{(\Alph*)}, leftmargin=2.2em, itemsep=2pt]
          \item 2
          \item 3
          \item \ans{5}
          \item 7
          \item 9
        \end{enumerate}

  \item Rainfall in millimetres is best described as
        \begin{enumerate}[label=\textbf{(\Alph*)}, leftmargin=2.2em, itemsep=2pt]
          \item a categorical variable, because each day is placed in a row.
          \item a discrete quantitative variable, because the recorded values are whole numbers.
          \item \ans{a continuous quantitative variable, because between any two rainfall amounts
                another amount is possible.}
          \item a continuous categorical variable.
          \item neither discrete nor continuous, because it was rounded when recorded.
        \end{enumerate}

  \item Suppose these 20 values are redrawn as a histogram with blocks 10 mm wide. Which
        statement is \textbf{true}?
        \begin{enumerate}[label=\textbf{(\Alph*)}, leftmargin=2.2em, itemsep=2pt]
          \item \ans{The histogram will show five bars of heights 5, 6, 4, 3, and 2.}
          \item The histogram will show 20 bars, one for each day.
          \item The histogram will let a reader recover that two days recorded exactly 13 mm.
          \item The histogram's bars must be drawn with gaps, as in a bar graph.
          \item The histogram cannot be drawn, because rainfall is continuous.
        \end{enumerate}

  \item A researcher redraws the same 20 values with blocks 2 mm wide and gets a jagged picture
        with many bars of height 0 or 1. The best conclusion is that
        \begin{enumerate}[label=\textbf{(\Alph*)}, leftmargin=2.2em, itemsep=2pt]
          \item the 2 mm version is wrong, because a histogram may not have empty blocks.
          \item the 10 mm version is wrong, because it hides individual values.
          \item \ans{both are correct; the interval width controls how much detail and how much
                overall shape the picture shows.}
          \item the data must have been recorded incorrectly.
          \item narrower blocks are always better, because they preserve more information.
        \end{enumerate}

\end{enumerate}

\vspace{0.14in}

\boxguard[16]
\begin{headlinebox}{goldbg}
\small
\textbf{Section II --- Free Response.} Show your reasoning. Every answer must be written
\emph{in the context of the problem} --- that is what earns credit on the AP exam.
\end{headlinebox}

\vspace{0.10in}

\begin{enumerate}[leftmargin=1.9em, itemsep=0.16in, resume]

  \item Use the rainfall record above.

        \vspace{6pt}
        \textbf{(a)} Classify rainfall as categorical or quantitative, and then as discrete or
        continuous. Justify from the variable itself, not from the display.\\[4pt]
        \ansline{Quantitative --- rainfall records a measured amount, not a category. It is}\\[1pt]
        \ansline{continuous: between 12 mm and 13 mm, 12.4 mm is perfectly possible.}

        \vspace{6pt}
        \textbf{(b)} Complete a frequency table for these 20 days using blocks 10 mm wide
        (0--9, 10--19, 20--29, 30--39, 40--49). Give the count for each block.\\[4pt]
        \ansline{0--9: 5 \quad 10--19: 6 \quad 20--29: 4 \quad 30--39: 3 \quad 40--49: 2.}\\[1pt]
        \ansline{The five counts add to 20, the number of rainy days recorded.}

        \vspace{6pt}
        \textbf{(c)} State one thing the stem plot shows that your table in part (b) does not,
        and one thing the table makes easier to see.\\[4pt]
        \ansline{The stem plot gives back every individual amount --- you can still read that}\\[1pt]
        \ansline{two days had exactly 13 mm. The table only says ``6 days between 10 and 19.''}\\[1pt]
        \ansline{The table makes it easier to see that light-rain days were the most common.}

        \vspace{6pt}
        \textbf{(d)} A colleague suggests displaying these 20 rainfall amounts as a bar graph
        with one bar per day. Explain why that is a poor choice here.\\[4pt]
        \ansline{Bar graphs display categories, and the days are not categories --- rainfall is}\\[1pt]
        \ansline{a quantitative variable on a number line. Twenty separate bars would show no}\\[1pt]
        \ansline{pattern in how much rain fell; grouping the amounts is what reveals the shape.}

\end{enumerate}

\vspace{0.12in}

\boxguard[12]
\begin{extensionbox}
\small
Find a histogram published in a news article or a report. Write down its interval width, then
say what the picture would probably hide if the width were doubled, and what it would probably
show if the width were halved.\\[4pt]
\ansline{Answers vary. Doubling the width typically hides a second peak or a gap;}\\[1pt]
\ansline{halving it typically exposes noise --- many bars of height 0 or 1.}
\end{extensionbox}

\vspace{0.12in}


\end{document}
